\documentclass[aspectratio=169,10pt]{beamer}

\usepackage[utf8]{inputenc}
\usepackage[T1]{fontenc}
\usepackage[spanish]{babel}
\usepackage{booktabs}
\usepackage{graphicx}
\usepackage{tikz}
\usetikzlibrary{arrows.meta,positioning}

\usetheme{Boadilla}
\usefonttheme{professionalfonts}
\definecolor{ReviewVerde}{RGB}{20,76,70}
\definecolor{ReviewTurquesa}{RGB}{0,126,120}
\definecolor{ReviewDorado}{RGB}{204,145,48}
\definecolor{ReviewCoral}{RGB}{180,82,58}
\definecolor{ReviewAzul}{RGB}{55,103,124}
\definecolor{ReviewOliva}{RGB}{126,123,54}
\definecolor{ReviewTinta}{RGB}{38,49,48}
\definecolor{ReviewClaro}{RGB}{239,245,243}
\definecolor{ReviewGris}{RGB}{112,122,120}

\setbeamercolor{structure}{fg=ReviewTurquesa}
\setbeamercolor{frametitle}{fg=white,bg=ReviewVerde}
\setbeamercolor{title}{fg=ReviewVerde}
\setbeamercolor{block title}{fg=white,bg=ReviewVerde}
\setbeamercolor{block body}{fg=ReviewTinta,bg=ReviewClaro}
\setbeamercolor{alerted text}{fg=ReviewDorado}
\setbeamertemplate{blocks}[default]
\setbeamerfont{frametitle}{size=\large,series=\bfseries}
\setbeamerfont{block title}{size=\normalsize,series=\bfseries}
\setbeamerfont{block body}{size=\small}
\newenvironment{reviewblock}[2]{%
  \begingroup
  \setbeamercolor{block title}{fg=white,bg=#1}%
  \begin{block}{#2}%
}{%
  \end{block}%
  \endgroup
}
\setbeamertemplate{navigation symbols}{}
\setbeamertemplate{itemize item}{\color{ReviewDorado}\small$\blacktriangleright$}
\setbeamertemplate{footline}{%
  \leavevmode\hbox{%
    \begin{beamercolorbox}[wd=.84\paperwidth,ht=2.6ex,dp=1ex,leftskip=.4cm]{author in head/foot}%
      \fontsize{8}{9}\selectfont Revisión sistemática | Hotspots de incidentes de seguridad
    \end{beamercolorbox}%
    \begin{beamercolorbox}[wd=.16\paperwidth,ht=2.6ex,dp=1ex,rightskip=.4cm plus1fil]{date in head/foot}%
      \fontsize{8}{9}\selectfont\hfill\insertframenumber/5
    \end{beamercolorbox}%
  }%
}

\title[Revisión de literatura]{Agrupamiento espacial y análisis espacio-temporal}
\subtitle{Revisión sistemática sobre hotspots de incidentes de seguridad}
\author{Ronaldo Ticona Jancco}
\institute{Universidad Nacional de San Antonio Abad del Cusco}
\date{Exposición breve | 5 minutos}

\begin{document}

% 1
\begin{frame}[plain]
  \begin{tikzpicture}[remember picture,overlay]
    \fill[ReviewVerde] (current page.north west) rectangle ([xshift=.34\paperwidth]current page.south west);
    \fill[ReviewDorado] ([xshift=.34\paperwidth]current page.north west) rectangle ([xshift=.355\paperwidth]current page.south west);
    \node[anchor=north west,text=white,font=\small\bfseries,align=left,text width=3.1cm]
      at ([xshift=.6cm,yshift=-1cm]current page.north west) {REVISIÓN\\DE LITERATURA};
    \node at ([xshift=.17\paperwidth]current page.west)
      {\includegraphics[height=3cm,keepaspectratio]{figures/Escudo_UNSAAC.png}};
    \node[anchor=south west,text=white,font=\small,align=left,text width=3.6cm]
      at ([xshift=.6cm,yshift=1cm]current page.south west) {Ronaldo Ticona Jancco\\UNSAAC | 2026};
  \end{tikzpicture}
  \vspace*{.75cm}
  \hspace*{.38\paperwidth}\begin{minipage}[t]{.57\paperwidth}
    \raggedright
    {\color{ReviewCoral}\small\bfseries INVESTIGACIÓN APLICADA}\par
    \vspace{.35cm}
    {\color{ReviewVerde}\LARGE\bfseries Hotspots de incidentes\\de seguridad}\par
    \vspace{.2cm}
    {\large Agrupamiento espacial y análisis espacio-temporal}\par
    \vspace{.7cm}
    {\color{ReviewAzul}\bfseries PROPÓSITO DE LA REVISIÓN}\par
    \vspace{.12cm}
    Revisar las técnicas empleadas, sus decisiones críticas y los criterios para evaluar la reproducibilidad del análisis de hotspots.\par
    \vspace{.65cm}
    {\color{ReviewOliva}\bfseries CASO QUE MOTIVA EL ESTUDIO}\par
    San Sebastián, Cusco\\
    Registros SIPCOP Municipal
  \end{minipage}
\end{frame}

% 2
\begin{frame}{Qué se revisó y cómo se seleccionó la literatura}
  \setbeamercolor{frametitle}{fg=white,bg=ReviewAzul}
  \begin{center}
  \begin{tikzpicture}[
    etapa/.style={draw=none,fill=ReviewClaro,rounded corners=1pt,minimum width=2.55cm,minimum height=1.05cm,align=center,font=\small\bfseries,text=ReviewTinta},
    paso/.style={-{Stealth[length=2mm]},thick,draw=ReviewCoral}]
    \node[etapa] (buscar) {1\\Buscar};
    \node[etapa,right=.55cm of buscar] (seleccionar) {2\\Seleccionar};
    \node[etapa,right=.55cm of seleccionar] (extraer) {3\\Extraer};
    \node[etapa,right=.55cm of extraer] (sintetizar) {4\\Sintetizar};
    \draw[paso] (buscar)--(seleccionar);
    \draw[paso] (seleccionar)--(extraer);
    \draw[paso] (extraer)--(sintetizar);
    \fill[ReviewCoral] (buscar.north west) rectangle ([xshift=.06cm]buscar.south west);
    \fill[ReviewAzul] (seleccionar.north west) rectangle ([xshift=.06cm]seleccionar.south west);
    \fill[ReviewTurquesa] (extraer.north west) rectangle ([xshift=.06cm]extraer.south west);
    \fill[ReviewOliva] (sintetizar.north west) rectangle ([xshift=.06cm]sintetizar.south west);
  \end{tikzpicture}
  \end{center}
  \vspace{.35cm}
  \begin{columns}[T,onlytextwidth]
    \column{.48\textwidth}
      \begin{reviewblock}{ReviewAzul}{Fuentes y temas}
        Búsquedas en IEEE Xplore, ScienceDirect y Scopus. Se abordaron patrones puntuales, hotspots, agrupamiento, análisis temporal y evaluación de artefactos.
      \end{reviewblock}
    \column{.48\textwidth}
      \begin{reviewblock}{ReviewTurquesa}{Selección y alcance}
        Se revisaron títulos, resúmenes y textos disponibles; se extrajeron métodos, datos, resultados y limitaciones. La síntesis es cualitativa: no se conservaron consultas literales ni conteos por etapa.
      \end{reviewblock}
  \end{columns}
\end{frame}

% 3
\begin{frame}{Cómo se organizó la evidencia encontrada}
  \setbeamercolor{frametitle}{fg=white,bg=ReviewCoral}
  \footnotesize
  \begin{columns}[T,onlytextwidth]
    \column{.48\textwidth}
      \begin{reviewblock}{ReviewCoral}{1. Patrones puntuales y hotspots}
        Densidad, kernel y proximidad describen concentración y recurrencia; dependen de la escala y del ancho de banda.
      \end{reviewblock}
      \begin{reviewblock}{ReviewAzul}{2. Agrupamiento espacial}
        DBSCAN identifica grupos y ruido; los métodos territoriales comparan unidades agregadas. No responden a la misma unidad de análisis.
      \end{reviewblock}
      \begin{reviewblock}{ReviewTurquesa}{3. Dimensión temporal}
        Ventanas, series y propagación representan aspectos distintos; la resolución disponible limita la interpretación.
      \end{reviewblock}
    \column{.48\textwidth}
      \begin{reviewblock}{ReviewOliva}{4. Calidad y trazabilidad}
        Coordenadas, duplicados, categorías y cobertura pueden alterar la forma y el tamaño de las concentraciones observadas.
      \end{reviewblock}
      \begin{reviewblock}{ReviewVerde}{5. Arquitectura y evaluación}
        La literatura orienta a describir componentes y responsabilidades, y a definir criterios observables de calidad y reproducibilidad.
      \end{reviewblock}
  \end{columns}
  \begin{center}
    \color{ReviewDorado}\bfseries Esta clasificación comparativa es el principal producto de síntesis de la revisión.
  \end{center}
\end{frame}

% 4
\begin{frame}{Hallazgos: qué condiciona la detección de hotspots}
  \setbeamercolor{frametitle}{fg=white,bg=ReviewTurquesa}
  \begin{columns}[T,onlytextwidth]
    \column{.52\textwidth}
      \begin{reviewblock}{ReviewCoral}{Los métodos son complementarios}
        KDE estima una superficie continua; DBSCAN forma grupos por densidad. La elección depende de la pregunta y la unidad de análisis.
      \end{reviewblock}
      \begin{reviewblock}{ReviewAzul}{Los resultados son sensibles}
        Escala, resolución de coordenadas, ancho de banda y parámetros $\varepsilon$ y MinPts pueden cambiar la concentración observada.
      \end{reviewblock}
    \column{.44\textwidth}
      \begin{reviewblock}{ReviewOliva}{La temporalidad tiene límites}
        Ventanas y series deben corresponder a la precisión temporal del registro; no se debe inferir una granularidad no disponible.
      \end{reviewblock}
      \begin{reviewblock}{ReviewVerde}{Grupo no significa significancia}
        Detectar un clúster no demuestra por sí solo un hotspot. Se requieren sensibilidad, referencia estadística y trazabilidad.
      \end{reviewblock}
  \end{columns}
  \vspace{.15cm}
  \begin{center}
    {\footnotesize Síntesis basada en Medina Ariza y Solymosi (2023), Cardone y Di Martino (2022), Cutid-Aguero et al. (2025), Changalasetty et al. (2023) y León Alarcón (2026).}
  \end{center}
\end{frame}

% 5
\begin{frame}{Qué obtuvo la revisión y qué aporta al caso local}
  \setbeamercolor{frametitle}{fg=white,bg=ReviewVerde}
  \begin{columns}[T,onlytextwidth]
    \column{.54\textwidth}
      \begin{reviewblock}{ReviewCoral}{Resultado de la revisión}
        Una síntesis de métodos, decisiones críticas y criterios de evaluación. No se identificó una arquitectura documentada y evaluada para integrar estas técnicas con los registros SIPCOP-M de San Sebastián.
      \end{reviewblock}
      \vspace{.15cm}
      \begin{reviewblock}{ReviewAzul}{Aplicación al caso}
        La revisión fundamenta un flujo que empieza con perfilado y validación de datos, selecciona técnicas según escala y pregunta, y termina con evaluación reproducible.
      \end{reviewblock}
    \column{.42\textwidth}
      \begin{reviewblock}{ReviewOliva}{Alcance y límite}
        La síntesis es cualitativa: no hubo metaanálisis ni conteos consolidados de selección. Sus resultados orientan el diseño local, no generalizaciones sobre todos los contextos urbanos.
      \end{reviewblock}
      \vspace{.15cm}
      \textbf{En una frase}\par
      El aporte no es un algoritmo nuevo, sino criterios para diseñar y evaluar un análisis adecuado a los datos disponibles.
  \end{columns}
\end{frame}

\end{document}